% =====================================================================
%  04-expedientes.tex - Expedientes, administracion penal y participantes
% =====================================================================

\section{Expedientes, administración penal y participantes}\label{sec:impl-expedientes}

El expediente es la unidad de aislamiento del despacho: delimita qué cuentas pueden consultar o modificar su información. Esta sección describe la pertenencia y el aislamiento, el perfil penal con revisiones administrativas y su cierre, el directorio de participantes y las identidades representadas con perfiles tipificados.

\subsection{Expedientes, pertenencia y aislamiento}\label{sec:impl-http-expedientes}

El dominio identifica cada expediente mediante \texttt{CaseId} y valida título y referencia con \texttt{Case\allowbreak Metadata}: ambos son obligatorios, se recortan espacios exteriores, se rechazan controles y se limitan a 200 y 100 caracteres, respectivamente. La referencia es una etiqueta libre, no un identificador único. PostgreSQL conserva una raíz inmutable con creador, fecha y metadatos originales; sus revisiones administrativas proporcionan título y referencia actuales. Las asignaciones regulan acceso y permanecen separadas de fichas de participantes, audiencias y plazos.

\begin{table}[H]
  \centering
  \caption{Contrato HTTP de expedientes y asignaciones. Los parámetros \texttt{<case>} y \texttt{<user>} son UUID.}
  \label{tab:impl-http-expedientes}
  \begin{tabularx}{\linewidth}{@{}L{6.2cm} X@{}}
    \toprule
    \textbf{Ruta} & \textbf{Responsabilidad} \\
    \midrule
    \texttt{POST /api/v1/cases} & Owner o Litigante crean, quedan asignados y confirman R1 pendiente sin etapa; \texttt{201}. \\
    \texttt{GET /api/v1/cases} & Lista únicamente metadatos visibles, con \texttt{limit} y \texttt{offset}. \\
    \texttt{GET /api/v1/cases/<case>} & Devuelve metadatos si el actor es Owner o miembro vigente. \\
    \texttt{PUT /api/v1/cases/\allowbreak <case>/members/<user>} & Owner asigna un usuario activo; operación idempotente, \texttt{204}. \\
    \texttt{DELETE /api/v1/cases/\allowbreak <case>/members/<user>} & Owner retira pertenencia; operación idempotente, \texttt{204}. \\
    \bottomrule
  \end{tabularx}
\end{table}

Owner ve todos los expedientes y administra asignaciones. Litigante puede crear un expediente, queda asignado automáticamente y pierde acceso si se retira esa asignación. La API básica conserva los cuatro campos de respuesta: UUID, título, referencia y creador. Cliente consulta esa proyección de sus expedientes asignados, sin perfil penal ni procedencia administrativa. Paralegal también dispone del índice y detalle del personal. Las lecturas básicas revalidan actor y alcance y confirman auditoría antes de responder. La consulta filtra pertenencia antes de paginar; el límite predeterminado es 50, admite de 1 a 100 y el desplazamiento es no negativo. El orden es por UUID, sin promesa de una instantánea común entre páginas. Un expediente ajeno responde igual que uno inexistente: \texttt{404 case\_not\_found}.

\subsection{Perfil penal y administración con revisiones}\label{sec:impl-http-administracion}

El módulo \rutarepo{crates/domain/src/case_administration/} separa los valores editables, el estado organizativo y las revisiones positivas. Un perfil presente exige NUC, autoridad emisora, carpeta judicial, órgano emisor y delitos registrados. La referencia interna conserva su significado de etiqueta libre y no sustituye esos identificadores. Los valores son declaraciones manuales: no verifican competencia institucional, identidad, naturaleza jurídica de un delito ni culpabilidad. El contrato y sus límites se documentan en \rutarepo{docs/case-administration-api.md} y \rutarepo{docs/adr/0022-audited-penal-case-administration.md}.

NUC y carpeta admiten hasta 100 escalares Unicode cada uno; sus autoridades, 200. La lista conserva entre uno y ocho delitos, de hasta 120 escalares por entrada, en el orden recibido. Rechaza duplicados literales después de recortar blancos exteriores. Información general admite hasta 1000 escalares y varias líneas: convierte CRLF a LF, rechaza CR aislado y los demás controles antes del recorte Unicode. Los otros textos rechazan todos los controles originales. Identificadores complementarios admite hasta 300 escalares; ambos campos opcionales vacíos quedan ausentes. Se conservan mayúsculas, acentos, forma Unicode y espacios interiores.

El canon \texttt{CADM1} comienza con cinco bytes ASCII y un byte de estado, cero para activo o uno para cerrado. Continúa con título y referencia, una etiqueta de presencia del perfil y, cuando existe, NUC, autoridad, carpeta, órgano, cantidad y textos ordenados de delitos y los dos opcionales. Cada texto lleva su longitud UTF-8 en cuatro bytes en orden de mayor a menor peso; cada opcional usa etiqueta cero o uno y, si existe, longitud y contenido. El máximo es 12\,717 bytes. SHA-256 cubre únicamente esos valores: expediente, revisión, autor y fecha se vinculan mediante claves y recurso auditado. El canon es propio del prototipo; no modifica el AAD ni la evidencia criptográfica documental y no equivale a una firma individual.

La raíz sin revisiones se presenta como R0, activa por compatibilidad, con perfil y procedencia ausentes. No existe una fila histórica cero. \texttt{Case\allowbreak Revision} representa únicamente valores de uno a $2^{32}-1$ y detecta agotamiento; la expectativa de escritura distingue expresamente una raíz sin historia de una revisión positiva. Toda alta básica por API confirma raíz, asignación del creador y R1 activa con perfil pendiente, junto con tres eventos. El alta penal completa añade el registro inicial de Investigación y confirma los cuatro eventos en la misma transacción. Ninguna operación devuelve una proyección reconstruida mediante un GET posterior al commit.

El registro inicial conserva número de etapa uno y referencia exacta a la revisión administrativa uno. Su huella, instante y autor UUID/correo se obtienen de esa revisión inmutable; una edición posterior no los sustituye. La fecha expresa el registro de la operación, no la fecha acreditada de un acto judicial. Completar un perfil básico o legado conserva la etapa sin registrar. La adopción explícita y los avances ordinarios se registran mediante el flujo separado de la \cref{sec:impl-http-etapas}.

\begin{table}[H]
  \centering
  \caption{Contrato HTTP de administración penal bajo \texttt{/api/v1}.}
  \label{tab:impl-http-administracion}
  \begin{tabularx}{\linewidth}{@{}L{5.4cm} X@{}}
    \toprule
    \textbf{Método y ruta} & \textbf{Resultado confirmado} \\
    \midrule
    \texttt{POST /penal-cases} & Alta completa activa, R1 y registro inicial; detalle con \texttt{201}. \\
    \texttt{GET /case-administrations} & Índice compacto autorizado del personal, filtros y cursor UUID. \\
    \texttt{GET /cases/<id>/}\allowbreak\texttt{administration} & Raíz, administración actual y registro inicial, cuando existe. \\
    \texttt{PUT /cases/<id>/}\allowbreak\texttt{administration} & Reemplazo de título, referencia y perfil con revisión esperada. \\
    \texttt{GET /cases/<id>/}\allowbreak\texttt{administration/history} & Revisiones positivas descendentes con cursor exclusivo. \\
    \texttt{PUT /cases/<id>/}\allowbreak\texttt{administrative-status} & Cambio exclusivo de estado con revisión esperada. \\
    \bottomrule
  \end{tabularx}
\end{table}

Owner consulta y administra todos los expedientes. Litigante asignado consulta, edita y cambia estado; Paralegal asignado consulta. Cliente recibe denegación en estas seis operaciones y conserva solamente la API básica. El adaptador vuelve a cargar cuenta activa, rol y pertenencia en la transacción; un UUID ajeno produce el mismo error que uno inexistente. Cada lectura, incluida una página vacía, confirma su evento antes de entregar datos.

El índice selecciona la cabeza antes de filtrar por estado, perfil pendiente o completo, título como subcadena literal y NUC o carpeta por igualdad exacta. Distingue mayúsculas y no interpreta comodines. Ordena UUID ascendentemente y usa \texttt{after\_id}; el historial ordena revisiones en sentido descendente y usa \texttt{before\_revision}. Ambos límites admiten de 1 a 100, con 50 por defecto. Una fila adicional determina si hay siguiente página; sin ella, el cursor es nulo. El índice omite autoridades, delitos, textos extensos, huellas y correos históricos. El detalle conserva origen, valores actuales y procedencia; la historia registra UUID y correo capturados sin consultar el correo actual de la cuenta.

La edición conserva estado y etapa y exige \texttt{expected\_revision}; un perfil completo ya no puede retirarse. El cambio exclusivo de estado lee y conserva los textos vigentes dentro de la transacción. Incluso repetir valores puede crear un sucesor. Una revisión obsoleta produce conflicto, sin reintento automático. NUC y carpeta son únicos por separado entre cabezas actuales de la instancia, incluidos los expedientes cerrados. Corregirlos libera los valores anteriores al confirmar y conserva su historia; el error no identifica el expediente que ocupa un valor. Esta regla de la instancia no acredita unicidad nacional.

El cierre administrativo impide editar el perfil, cargar, añadir versiones, clasificar y sellar documentos, además de crear, editar, archivar o reactivar participantes y registrar adopciones o transiciones de etapa. Mantiene consultas, historias, verificación y exportación autorizadas; Owner conserva asignación y revocación. Reactivar el expediente no reactiva fichas archivadas ni modifica la etapa. Las mutaciones comprueban el cierre después de autorizar y resolver el recurso exacto, también al confirmar una preparación criptográfica anterior. Si el cierre confirma primero, esa preparación no produce nueva fila ni evento de éxito. El estado organizativo no concluye ni reabre un proceso judicial.

Los cuerpos JSON de alta, edición y estado tienen límite de 64~KiB sobre el cuerpo completo. Rechazan campos desconocidos o repetidos, también dentro del perfil, y no aceptan autor, fecha, huella o etapa suministrados por el cliente. Las entradas mal formadas producen \texttt{400}; un cuerpo excesivo, \texttt{413}; los valores fuera del contrato, \texttt{422}. Revisión obsoleta o agotada, cierre, identificador ocupado y retirada de perfil completo producen \texttt{409}. La inconsistencia persistida se informa como error interno sin detalles de almacenamiento. Una respuesta perdida puede suceder después del commit: el cliente debe conservar los datos y reconciliar el estado autorizado antes de decidir otro envío.

\subsection{Directorio de participantes por expediente}\label{sec:impl-http-participantes}

El directorio conserva el registro manual de fichas con UUID propio y asociación inmutable a un expediente. Una ficha puede representar a una persona o autoridad mencionada por el despacho; no crea una cuenta ni una asignación de acceso. En el registro manual se permiten homónimos, incluso con el mismo rol escrito. La extensión tipificada vincula identidades y perfiles mediante las reglas de la \cref{sec:impl-participantes-tipificados}. La cuenta autenticada que crea o modifica permanece separada de la persona representada. El contrato manual se documenta en \rutarepo{docs/adr/0021-audited-case-participants.md}.

El valor de dominio \texttt{Participant\allowbreak Values} exige nombre y rol manual, de hasta 200 y 80 caracteres Unicode, respectivamente. Organización y estatus jurídico son opcionales, con máximos de 200 y 160. Rechaza caracteres de control en la entrada original antes de recortar espacios Unicode exteriores; una opción vacía queda ausente. Conserva espacios interiores, mayúsculas, acentos, puntuación y forma de normalización Unicode. El estado del directorio es \texttt{active} o \texttt{archived}: organiza la vista, sin cambiar el estatus jurídico ni el del expediente.

Cada ficha comienza en revisión uno y estado activo. El tipo \texttt{Participant\allowbreak Revision} acepta de uno a $2^{32}-1$; no existe una revisión cero ni se inventa procedencia inicial. El reemplazo completo exige la revisión actual esperada y añade un sucesor, incluso si los valores coinciden. Archivar o reactivar usa un comando separado: lee la revisión actual y cambia únicamente el estado dentro de la misma transacción. Así conserva una edición de texto ya confirmada; si la revisión esperada quedó obsoleta, devuelve conflicto sin aplicar el cambio.

El canon \texttt{PART1} codifica primero cinco bytes ASCII, después nombre y rol como longitud UTF-8 de cuatro bytes en orden de mayor a menor peso y su contenido. Cada opcional usa un byte cero si está ausente o uno seguido de longitud y contenido; el estado termina con cero para activo o uno para archivado. La longitud máxima es 2\,584 bytes. SHA-256 cubre estos valores, mientras el recurso de auditoría vincula expediente, ficha, revisión y digest. Esta huella no equivale a una firma individual ni acredita identidad jurídica.

\texttt{Participant\allowbreak Service} autentica mediante el puerto de identidad; \texttt{Participant\allowbreak Store} delimita la persistencia. Owner administra todos los directorios, Litigante asignado administra y Paralegal asignado consulta. Cliente no tiene acceso, aunque pueda leer los metadatos básicos de su expediente. El adaptador PostgreSQL vuelve a consultar cuenta activa, rol y pertenencia dentro de la transacción auditada. También Owner debe proporcionar el expediente exacto de la ficha. Una ficha ausente o ajena devuelve el mismo error, antes de comprobar la revisión esperada.

La raíz, la primera revisión y el evento de creación se confirman juntos. Cada sucesor conserva valores canónicos, digest, instante UTC y UUID y correo del mismo actor recargado durante la operación. El historial no reconstruye autores mediante sus perfiles actuales. La fecha conserva el instante de la operación; revisión y secuencia de auditoría ordenan las confirmaciones sin prometer marcas temporales estrictamente crecientes. Fallos de inserción, auditoría o commit revierten la mutación completa. La respuesta representa el estado confirmado por esa operación, sin una consulta posterior.

{\emergencystretch=3em La \cref{tab:impl-http-participantes} usa el prefijo \rutarepo{/api/v1/cases/<case>/participants}. Las mutaciones manuales conservan su respuesta plana: expediente, UUID, revisión, los cinco valores, \texttt{values\_digest}, \texttt{changed\_at} y \texttt{changed\_by}, este último con UUID y correo capturado. Lecturas, historia y estado admiten la unión de fichas manuales y tipificadas; el listado presenta una proyección compacta.\par}

\begin{table}[H]
  \centering
  \caption{Contrato HTTP del directorio de participantes.}
  \label{tab:impl-http-participantes}
  \begin{tabularx}{\linewidth}{@{}L{4.3cm} X@{}}
    \toprule
    \textbf{Método y sufijo} & \textbf{Semántica} \\
    \midrule
    \texttt{POST} (sin sufijo) & Recibe cuatro textos, fija estado activo y devuelve la revisión uno con \texttt{201}. \\
    \texttt{PUT /<id>} & Reemplaza textos y estado con \texttt{expected\_revision}; devuelve su revisión confirmada. \\
    \texttt{PUT /<id>/}\allowbreak\texttt{directory-status} & Cambia solo estado con revisión esperada, preservando textos actuales. \\
    \texttt{GET} (sin sufijo) & Consulta cabezas actuales con filtros y cursor exclusivo de UUID. \\
    \texttt{GET /<id>} & Devuelve la revisión actual de la ficha autorizada. \\
    \texttt{GET /<id>/history} & Devuelve revisiones históricas en orden descendente con cursor exclusivo. \\
    \bottomrule
  \end{tabularx}
\end{table}

El listado toma primero la última revisión de cada ficha manual o tipificada y después combina nombre por subcadena literal y estado. El filtro de rol manual por igualdad exacta solo devuelve fichas pendientes de tipificar; tipo y perfil permiten consultar las variantes actuales. Nombre y rol distinguen mayúsculas y formas Unicode; los signos de comodín no expanden coincidencias. El estado predeterminado es activo y puede solicitarse archivado o ambos. La paginación ordena UUID de forma ascendente y utiliza \texttt{after\_id}; historia ordena revisiones de forma descendente con \texttt{before\_revision}. Ambos límites van de 1 a 100, con 50 por defecto. Una fila adicional determina \texttt{has\_more} y el cursor siguiente; sin ella, el cursor es nulo. No se calcula un total global ni se promete una instantánea compartida entre páginas.

Cada lectura confirma su evento antes de devolver datos, incluso con una página vacía. Los eventos del contrato manual distinguen creación, reemplazo, cambio exclusivo de estado, listado, detalle e historia, sin copiar nombres ni etiquetas jurídicas al recurso de auditoría. Las escrituras manuales rechazan campos desconocidos o repetidos y limitan el JSON completo a 8~KiB, incluidos bytes posteriores al objeto. Cero en una revisión produce \texttt{422}; tipos incorrectos, UUID inválido o JSON mal formado producen \texttt{400}; un cuerpo excesivo, \texttt{413}. Los conflictos distinguen revisión obsoleta y contador agotado. La corrupción persistida se informa como error interno sin exponer detalles.

La migración incorpora un directorio vacío a expedientes existentes; no convierte asignaciones de usuarios en participantes. PostgreSQL exige revisión inicial activa, secuencia contigua, canon y huella válidos, referencias al expediente y actor, y una relación diferida que impide confirmar una raíz sin revisión uno. El rol operativo puede leer e insertar, pero no actualizar, eliminar, truncar ni poseer los objetos protegidos. El arranque comprueba esquema, privilegios e inventario. Las funciones resuelven sus referencias en el esquema de instalación, también durante restauración con \texttt{search\_path} vacío. El respaldo conserva raíces, revisiones, procedencia y auditoría junto con documentos y evidencia. Estas medidas no detectan por sí solas una restauración administrativa antigua y coherente.

\subsubsection{Límites de la representación jurídica}

El rol de una ficha pendiente es una descripción manual, no una regla de autorización. El artículo 105 del CNPP distingue sujetos del procedimiento y quienes tienen calidad de parte \cite{cnpp_consolidado2025}. Los perfiles tipificados agregan datos específicos y la declaración interna comprueba una firma real, pero ninguno acredita nombramiento o representación. Los artículos 6 y 7 del acuerdo FIREL del OAJ de 2026 regulan firma en su sistema electrónico y reconocimiento procesal para actuar por otra persona \cite{oaj_firel2026}; poseer un certificado no sustituye ese reconocimiento. La separación técnica entre directorio, acceso y firma conserva esa distinción.

El acuerdo de expediente electrónico AG-POAJ-006/2026 abroga el Acuerdo General 12/2020 \cite{oaj_expediente2026}. El acuerdo FIREL AG-POAJ-016/2026 deroga las disposiciones de emisión, administración, renovación, revocación y uso del acuerdo conjunto 1/2013 en el ámbito de competencia del OAJ, sin declarar la abrogación íntegra de ese antecedente \cite{oaj_firel2026}. Estas fuentes describen los procedimientos judiciales a los que se refieren; no convierten el registro interno de una ficha en una actuación judicial.

El caso de uso del directorio permanece parcial respecto de los criterios jurídicos: datos por tipo, identidad estable local, revisión de coincidencias y firma interna no equivalen a identidad civil comprobada, FIREL o representación acreditada. El estado abierto del expediente es administrativo, no una certificación de condición procesal. El alcance se conserva en \rutarepo{docs/product-completion.md}; las políticas de datos personales, retención e interoperabilidad judicial requieren validación propia.

\subsection{Identidades representadas y participantes tipificados}\label{sec:impl-participantes-tipificados}

Junto a las fichas manuales de la \cref{sec:impl-http-participantes}, el directorio registra identidades locales del expediente con revisión propia. \texttt{Case\allowbreak SubjectId} identifica a la persona u órgano; \texttt{ParticipantId} identifica la ficha y su rol; \texttt{UserId} identifica a la cuenta autora. Crear una identidad junto con su primera ficha es atómico. Reutilizarla fija una revisión exacta; modificarla después no reescribe la identidad vinculada a fichas o firmas anteriores. La consulta de identidad actual constituye una operación explícita distinta.

\begin{table}[H]
  \centering
  \caption{Datos propios de los once perfiles del directorio.}
  \label{tab:impl-perfiles-tipificados}
  \begin{tabularx}{\linewidth}{@{}L{3.8cm} X@{}}
    \toprule
    \textbf{Perfil} & \textbf{Datos declarados y condiciones particulares} \\
    \midrule
    Imputado & Situación de libertad conocida o motivo de desconocimiento. \\
    Víctima u ofendido & Contacto y protección mediante estados explícitos y soportes cuando están documentados. \\
    Defensor & Cédula y emisor, modalidad pública o particular; declaración firmada obligatoria. \\
    Ministerio Público & Identificador de fiscalía, unidad y cédula, con desconocimiento motivado admisible. \\
    Asesor jurídico & Institución o despacho y cédula conocida o desconocida con motivo. \\
    Juez de control & Órgano jurisdiccional y declaración firmada obligatoria. \\
    Tribunal de enjuiciamiento & Identidad institucional, distrito y composición unipersonal o colegiada. \\
    Perito & Especialidad y cédula, con desconocimiento motivado admisible. \\
    Policía & Corporación y unidad declaradas o desconocidas con motivo. \\
    Supervisión cautelar & Autoridad y unidad declaradas o desconocidas con motivo. \\
    Otro & Etiqueta y descripción; admite persona física u órgano compatible. \\
    \bottomrule
  \end{tabularx}
\end{table}

Organización y estatus jurídico permanecen como valores organizativos opcionales. La identidad y el rol conservan soportes cifrados con UUID documental, versión, SHA-256 y localizador. El conjunto de una operación, incluidas decisiones de candidatos, admite como máximo dos versiones distintas, 16~MiB cada una y 32~MiB en total. La preparación verifica integridad, evidencia capturada y admisión de formato mediante el worker descrito en la \cref{sec:impl-formatos}. Seleccionar un soporte no modifica su contenido; una carga confirmada permanece si la preparación posterior es rechazada.

\subsubsection{Canon, revisión y confirmación}

\texttt{SUBJ1} codifica la identidad y admite hasta 5676 bytes. \texttt{PART2} codifica el rol y su referencia exacta de identidad, hasta 8380 bytes. Las revisiones manuales se codifican con \texttt{PART1}. Los enteros tienen tamaño y orden definidos, las variantes usan etiquetas y los textos conservan UTF-8. Los codecs rechazan representaciones no canónicas, campos incompatibles y datos posteriores al valor completo. SHA-256 cubre cada codificación; una revisión de ficha tipificada no incorpora su propia firma al canon, evitando una dependencia circular.

La revisión de candidatos combina nombre conocido, identificador, certificado y pareja de digest documental y localizador. Una etiqueta de persona no identificada no equivale a un nombre conocido. Las decisiones de identidad distinta deben cubrir exactamente los candidatos devueltos, con motivo y soporte. La capacidad máxima es de 16; excederla rechaza la preparación y no recomienda alterar datos para eludirla. La huella del directorio se obtiene recorriendo sus cabezas ordenadas con memoria acotada. Una edición concurrente fuera de la página visible también invalida esa revisión.

{\raggedright Bajo \rutarepo{/api/v1/cases/<case>/participants/proposals}, las rutas \texttt{review}, \texttt{prepare} y \texttt{commit} separan revisión, preparación y confirmación. La primera propone UUID sin reservarlos; la segunda devuelve expectativas, valores y evidencia a firmar cuando aplica; la tercera confirma la operación completa. Sobre el prefijo del directorio, \texttt{/<id>/revisions/<revision>} devuelve el detalle exacto de una revisión manual o tipificada, y el sufijo adicional \texttt{/credential}, la evidencia pública asociada a esa revisión; la historia también reúne ambas familias. El recurso \texttt{subjects} ofrece índice compacto, detalle actual, revisión exacta, historia y reemplazo revisado de identidad. El contrato detallado se conserva en \rutarepo{docs/typed-participants-api.md}.\par}

Las rutas tipificadas limitan el JSON completo a 256~KiB y exigen fin de entrada; las escrituras manuales se limitan a 8~KiB. Los índices parten de las cabezas actuales: el listado de fichas aplica sus filtros sobre la cabeza de la unión manual y tipificada, como describe la \cref{sec:impl-http-participantes}. Identificadores, contacto y material de credenciales se excluyen de los índices; su detalle requiere una consulta explícita con la política de lectura vigente.

La persistencia conserva raíces y revisiones de identidad, revisiones tipificadas, decisiones de candidatos y evidencia pública. Un trigger diferido exige exactamente una revisión inicial manual o tipificada para cada ficha. Bajo la transacción auditada se revalidan cuenta activa, rol, pertenencia, estado abierto, revisiones esperadas, unicidad de identidad y rol, huella del directorio y soportes completos. La sesión también se comprueba antes de ingresar a persistencia. Un fallo revierte valores, decisiones y auditoría juntos. El cierre administrativo impide nuevas mutaciones y conserva lecturas y evidencias autorizadas (\cref{sec:impl-http-administracion}). Los permisos por rol coinciden con los del directorio manual (\cref{sec:impl-http-participantes}).

\subsubsection{Declaración personal y confianza publicada}\label{sec:impl-declaracion-participante}

El perfil \texttt{internal\_demo\_v1} valida una cadena raíz-hoja acotada; la hoja aportada por el operador no elige su raíz de confianza. Exige RSA-3072, usos de clave de firma adecuados, extensiones admitidas y CRL directa completa vigente. No realiza consultas AIA, OCSP ni descubrimiento de intermediarios. Certificado y raíz admiten 16~KiB cada uno; la CRL admite 1~MiB y 10\,000 entradas. Material mal formado, revocado, vencido o fuera del perfil se rechaza. Estas condiciones son propias de la demostración interna, no un validador general de FIREL.

La declaración \texttt{PCRED1} mide exactamente 218 bytes. Fija despliegue y raíz, expediente, identidad, operación y revisiones de identidad, ficha y revisiones de ficha, tipo de rol, digests de valores y huella DER de la hoja. Se exportan esos bytes y un recibo legible antes de firmar. La firma separada tiene exactamente 384 bytes y usa RSASSA-PKCS1-v1\_5 con SHA-256, sin aplicar un segundo hash al digest ya preparado. El titular conserva la clave privada fuera del servidor; no se admiten contraseñas de clave, PFX ni PKCS\#12. La clave compartida de sellado documental no sustituye esta firma personal.

El administrador publica raíz y CRL mediante revisiones inmutables de confianza en PostgreSQL. Se conserva el usuario administrativo de sesión como procedencia, sin inventarle una cuenta de Qadra. La raíz del despliegue permanece fija; el número de CRL crece estrictamente y su fecha de emisión no retrocede. Generar o editar el índice local de la CA no actualiza por sí mismo la confianza publicada. El rol operativo solo lee estos materiales. Después de verificar criptografía fuera del bloqueo de auditoría, la confirmación exige la misma cabeza de confianza y captura el reloj tras esperar el bloqueo; el instante debe permanecer dentro del intervalo admitido de raíz, hoja y CRL.

La evidencia histórica conserva declaración, firma, certificados, CRL, revisión de confianza, referencias exactas y momento comprobado. Ese momento no se declara fecha original de firma. Una CRL posterior no reescribe el resultado histórico. Archivar o reactivar conserva la referencia de la credencial anterior, sin afirmar que la firma cubra el nuevo estado organizativo; editar valores que requieren credencial exige una nueva firma. El respaldo debe conservar todo el conjunto. Un rollback administrativo coherente de toda la base sigue fuera de la garantía sin anclaje externo.

\subsubsection{Interacción y conciliación en Qadra}

El alta principal presenta el formulario tipificado; registrar una ficha pendiente mantiene el flujo manual explícito. El operador selecciona una identidad existente o registra una nueva, revisa candidatos, prepara la declaración y confirma. El certificado se libera al cambiar identidad, expediente o sesión; cambiar valores invalida revisión, declaración y firma, pero puede conservar el certificado de la misma identidad. Borradores y archivos permanecen solo en memoria (\cref{sec:impl-qadra-borradores}). Las acciones vecinas esperan el trabajo real de consultas o preparación, sin lanzar lecturas por cada fila del índice o de la historia.

Ante una respuesta incierta se conserva el envío exacto. Para una operación sin firma, la revisión propuesta se concilia mediante origen y digest \texttt{PTXN1}; para una firmada, mediante declaración, firma, certificado y referencias originales, sin aceptar como éxito una coincidencia de texto actual. La edición separada de identidad permite consultar el sucesor esperado y decidir explícitamente antes de reemplazar otra vez; si esa revisión aún no existe, conserva el borrador. Conflictos, denegaciones y respuestas tardías siguen el tratamiento común de las \cref{sec:impl-qadra-borradores,sec:impl-qadra-conciliacion}, que conserva los valores locales o retira el detalle según el caso; ningún camino reenvía automáticamente ni genera nuevos UUID para ocultar la incertidumbre.

Ni el formulario ni la declaración firmada amplían los límites jurídicos del directorio reunidos en la \cref{sec:impl-participantes-limites}.

